\begin{figure}[htbp]
\centering
\begin{tikzpicture}[causal,x=1.5cm,y=1.2cm]
 \node[latent] (ch09f8-h) at (.75,1) {$H$};
 \node[latent] (ch09f8-j) at (3,1) {$J$};
 \node[vtx] (ch09f8-a) at (0,0) {$A$};
 \node[vtx] (ch09f8-b) at (1.5,0) {$B$};
 \node[vtx] (ch09f8-c) at (2.25,0) {$C$};
 \node[vtx] (ch09f8-d) at (3.75,0) {$D$};
 \draw[ed] (ch09f8-h)--node[left] {$\alpha$}(ch09f8-a);
 \draw[ed] (ch09f8-h)--node[below left] {$\beta$}(ch09f8-b);
 \draw[ed] (ch09f8-h)--node[above] {$\eta$}(ch09f8-j);
 \draw[ed] (ch09f8-j)--node[below right] {$\gamma$}(ch09f8-c);
 \draw[ed] (ch09f8-j)--node[right] {$\delta$}(ch09f8-d);
\end{tikzpicture}
\caption{如果该图对应一个线性结构因果模型，其导出的分布将满足四元约束(9.12)—(9.14)。}
\end{figure}
